\section{Circuits: from components to algorithms}
\label{sec:circuits}

The tools of the last sections answer the first two questions of \cref{sec:intro}: what is represented, and
which components matter. A \emph{circuit} is an answer to the third: how the components work together.

\subsection{What counts as a circuit}

By \cref{prop:sum}, the input of every component is a sum of the outputs of earlier components. We can
therefore draw the model as a graph. The nodes are the embeddings, the attention heads and the MLPs. There is
an edge from every component to every later component and to the output, which stands for the term that the
earlier one contributes to the input of the later one. This graph has the same structure for every model of
the same architecture. What differs is which edges carry something that matters.

\begin{definition}[Circuit]\label{def:circuit}
A \emph{circuit} for a task is a subgraph of this graph together with a description of what each of its nodes
computes on the task.
\end{definition}

A drawing of nodes and arrows is easy to produce, so the definition is only useful together with tests.
\citet{wang2023interpretability} proposed three. In all of them, components outside the set under study are
replaced by their mean over a reference distribution of prompts (mean ablation, \cref{sec:interventions}).
\begin{itemize}
\item \textbf{Faithfulness.} The circuit alone, with everything else ablated, performs about as well as the
full model.
\item \textbf{Completeness.} The circuit contains all components that take part in the behaviour. The test
removes the same subset of nodes from the circuit and from the full model and asks that the two still perform
alike, for every subset. It detects components outside the circuit that could step in.
\item \textbf{Minimality.} Every node of the circuit is needed: for each node there is a situation in which
removing it changes the performance noticeably.
\end{itemize}

\subsection{Reading an algorithm from the weights}
\label{sec:pathexp}

In sufficiently small models a circuit can be read from the weights, without running the model at all. The
cleanest case is a transformer with one layer and no MLP, studied by \citet{elhage2021mathematical}. To keep the
formulas short we also drop $\LN$ and the positional embedding. The model is then
\begin{equation}\label{eq:onelayer}
x_i^{(0)}=W_E[t_i],\qquad
x_i^{(1)}=x_i^{(0)}+\sum_{h=0}^{H-1}\sum_{j\le i}A^{0,h}_{ij}\,W_{OV}^{0,h}\,x_j^{(0)},\qquad
\logit_i(t)=W_U[t]^\top x_i^{(1)},
\end{equation}
where $W_{OV}^{0,h}$ is the matrix $W_{OV}$ of head $h$ in the single layer $0$.

\begin{proposition}[Path expansion of a one-layer attention-only transformer]\label{prop:path}
In the model \eqref{eq:onelayer},
\begin{equation}\label{eq:path}
\logit_i(t) \;=\; \underbrace{W_U[t]^\top W_E[t_i]}_{\text{direct path}}
\;+\;\sum_{h=0}^{H-1}\sum_{j\le i}A^{0,h}_{ij}\;\underbrace{W_U[t]^\top W_{OV}^{0,h}\,W_E[t_j]}_{\text{path through head }h} .
\end{equation}
\end{proposition}
\begin{proof}
We insert the stream $x_i^{(1)}$ of \eqref{eq:onelayer} into the read-out.
\begin{align*}
\logit_i(t) &= W_U[t]^\top x_i^{(0)}+\sum_{h=0}^{H-1}\sum_{j\le i}A^{0,h}_{ij}\,W_U[t]^\top W_{OV}^{0,h}\,x_j^{(0)}
&& \by{by \eqref{eq:onelayer}, linearity}\\
&= W_U[t]^\top W_E[t_i]+\sum_{h=0}^{H-1}\sum_{j\le i}A^{0,h}_{ij}\,W_U[t]^\top W_{OV}^{0,h}\,W_E[t_j].
&& \by{since $x_j^{(0)}=W_E[t_j]$}\qedhere
\end{align*}
\end{proof}

Every term of \eqref{eq:path} is an entry of a table with one row and one column per token, and each table has
a plain meaning.
\begin{itemize}
\item $W_U[t]^\top W_E[t_i]$ says how much the current token $t_i$ promotes the next token $t$ by itself. This
table approximates the statistics of pairs of adjacent tokens.
\item $W_U[t]^\top W_{OV}^{0,h}W_E[t_j]$ says how much token $t$ is promoted \emph{if} head $h$ attends to an
earlier token $t_j$.
\item Whether it attends there is decided by a third table. Without $\LN$, the score in \eqref{eq:pattern} is
proportional to $W_E[t_i]^\top W_{QK}^{0,h}W_E[t_j]$.
\end{itemize}
A one-layer attention-only transformer is therefore a sum of look-up tables of the form ``when the current
token is $t_i$ and an earlier token is $t_j$, raise or lower the probability of $t$''. The three tables can be
computed from the weights and inspected. The explanation then covers all inputs, and not only the ones that
were tried.

With two layers, the input of a head in the second layer contains the outputs of the heads in the first, and
\eqref{eq:path} acquires terms in which the tables of two heads are multiplied. This is called
\emph{composition}. It comes in three kinds, according to whether the earlier head influences the later head's
query (where it looks from), its key (what it looks for) or its value (what it moves). Composition is where
real algorithms begin.

\subsection{Induction heads}
\label{sec:induction}

The first algorithm found in this way is also the most important one known. Text is full of repetitions. If the
sequence \tok{[A][B]} has occurred and the token \tok{[A]} occurs again, then \tok{[B]} is a good guess for the
next token. Two-layer attention-only transformers implement this rule with two heads
\citep{elhage2021mathematical}, as shown in \cref{fig:induction}.

\begin{enumerate}
\item A \emph{previous-token head} in the first layer attends from every position to the position before it.
At the position of \tok{[B]} it writes into the stream the information ``the token before me was \tok{[A]}''.
\item An \emph{induction head} in the second layer sits at the second occurrence of \tok{[A]}. Its query asks
for positions whose stream says ``the token before me was \tok{[A]}''. The keys it compares with contain the
output of the previous-token head, which is a case of key composition. The head therefore attends to \tok{[B]}.
\item The matrix $W_{OV}$ of the induction head copies: it raises the logit of the token it attends to. The
model predicts \tok{[B]}.
\end{enumerate}

\begin{figure}[t]
\centering
\begin{tikzpicture}[>=Latex, font=\small, tokn/.style={draw=tutgrey, fill=tutpaper, minimum height=6mm, minimum width=9mm}]
\node (l) at (0,0) {$\cdots$};
\node[tokn] (a1) at (1.3,0) {\tok{[A]}};
\node[tokn] (b1) at (2.7,0) {\tok{[B]}};
\node (m) at (4.6,0) {$\cdots$};
\node[tokn] (a2) at (6.5,0) {\tok{[A]}};
\node[tokn, dashed, fill=white] (b2) at (7.9,0) {\tok{[B]}};
\draw[->, line width=1pt, tutblue] (a1.south) to[out=-60,in=-120] node[below, text=black, align=center]
  {(1) previous-token head, layer 0:\\ ``the token before me was \tok{[A]}''} (b1.south);
\draw[->, line width=1pt, tutred] (b1.north) to[out=50,in=130] node[above, text=black, align=center]
  {(2) induction head, layer 1: attends to the token that follows an earlier \tok{[A]}} (a2.north);
\draw[->, line width=1pt, tutred] (a2.east) -- node[below=3mm, text=black, align=center, xshift=14mm] {(3) copies \tok{[B]}\\ to the prediction} (b2.west);
\end{tikzpicture}
\caption{The induction circuit. Arrows show the flow of information from a source position to a destination
position. The dashed box is the predicted next token.}
\label{fig:induction}
\end{figure}

Induction heads are not a curiosity of toy models. \citet{olsson2022context} found them in models of many
sizes and showed that they form during a short window early in training, at the same time as the model's
ability to use its context improves sharply. They are a mechanism of \emph{in-context learning}: the model picks
up a pattern from the prompt and continues it. Heads 5.5 and 6.9 of GPT-2 small are induction heads, and we have
already met them in \cref{tab:heads}.

\subsection{The IOI circuit}
\label{sec:ioi}

We can now assemble the evidence on the running example into the circuit of \citet{wang2023interpretability},
shown in \cref{fig:ioi}. As in their paper, we work backwards from the output.

\paragraph{Who writes the answer?} The name mover heads 9.9, 9.6 and 10.0. At the last position they attend to
the indirect object (\cref{sec:attention}) and write in the direction that raises its logit (\cref{fig:dla}).
Their $W_{OV}$ copies the name they attend to.

\paragraph{Why do they attend to the right name?} Both names are candidates. The name movers avoid the subject
because of the S-inhibition heads 7.3, 7.9, 8.6 and 8.10. These heads sit at the last position and attend to
S2. They write a signal into the last position that changes the queries of the name movers, so that positions
holding the subject receive less attention. This is the transfer from S2 to the last position in layers 7 and 8
that we saw in \cref{fig:residpatch}, and it explains why these heads have a large patching effect and a small
direct effect in \cref{tab:heads}.

\paragraph{How do the S-inhibition heads know which name is repeated?} Two groups of heads at S2 detect the
repetition. The duplicate token heads 0.1 and 3.0 attend from S2 to the earlier occurrence S1 of the same
token. The induction heads 5.5 and 6.9 attend from S2 to the token \emph{after} S1, using the output of the
previous-token heads 2.2 and 4.11 exactly as in \cref{sec:induction}. Both groups write to the stream at S2 that
this token has occurred before, and the S-inhibition heads read it.

\begin{figure}[t]
\centering
\begin{tikzpicture}[>=Latex, font=\footnotesize,
  tokn/.style={draw=tutgrey, fill=tutpaper, minimum height=5.5mm, inner sep=3pt},
  cls/.style={draw=tutgrey, rounded corners=2pt, align=center, inner sep=3.5pt, fill=#1}, cls/.default=tutblue!10]
\node[tokn] (t1) at (0,0) {\tok{When}};
\node[tokn] (io) at (1.4,0) {\tok{Mary}};
\node[tokn] (t3) at (2.6,0) {\tok{and}};
\node[tokn] (s1) at (3.8,0) {\tok{John}};
\node[tokn] (s1p) at (5.1,0) {\tok{went}};
\node (d1) at (6.3,0) {$\cdots$};
\node[tokn] (s2) at (7.7,0) {\tok{John}};
\node (d2) at (9.4,0) {$\cdots$};
\node[tokn] (end) at (11.2,0) {\tok{to}};
\node[below=1pt of io, tutgrey] {IO};
\node[below=1pt of s1, tutgrey] {S1};
\node[below=1pt of s2, tutgrey] {S2};
\node[below=1pt of end, tutgrey] {last position};
\node[cls] (prev) at (5.1,2.2) {previous token\\ 2.2, 4.11};
\node[cls] (dup) at (7.7,1.05) {duplicate token\\ 0.1, 3.0};
\node[cls] (ind) at (7.7,2.2) {induction\\ 5.5, 6.9};
\node[cls] (sinh) at (11.2,3.3) {S-inhibition\\ 7.3, 7.9, 8.6, 8.10};
\node[cls=tutred!12] (nm) at (9.3,5.3) {name movers\\ 9.9, 9.6, 10.0};
\node[cls=tutblue!22] (neg) at (13.0,5.3) {negative name movers\\ 10.7, 11.10};
\node (out) at (11.2,6.9) {logit difference in favour of \tok{Mary}};
\node at (4.7,5.6) {copy the attended name};
\draw[->, tutgrey] (s1.north) -- (prev.south west);
\draw[->, tutgrey] (s1.north east) -- (dup.west);
\draw[->, tutgrey] (prev.east) -- (ind.west);
\draw[->, tutgrey] (dup.east) to[out=0,in=200] (sinh.south west);
\draw[->, tutgrey] (ind.east) to[out=0,in=180] (sinh.west);
\draw[->, tutgrey] (sinh.north) to[out=120,in=-50] node[right=2pt, pos=0.5, text=black] {avoid S} (nm.south east);
\draw[->, tutred, line width=0.9pt] (io.north) to[out=80,in=180] (nm.west);
\draw[->, tutred, line width=0.9pt] (nm.north) -- node[left=2pt, text=black] {$+$} (out.south west);
\draw[->, tutblue, line width=0.9pt] (neg.north) -- node[right=2pt, text=black] {$-$} (out.south east);
\end{tikzpicture}
\caption{The circuit for indirect object identification in GPT-2 small, simplified from
\citet{wang2023interpretability}. Each box is a class of attention heads, drawn above the position at which the
heads write. Arrows show the flow of information. The backup name movers are omitted.}
\label{fig:ioi}
\end{figure}

\paragraph{The algorithm.} In words, the model does three things. It finds the name that occurs twice. It
tells the last position to ignore that name. It copies the name that is left. Two further classes complicate
the picture. The negative name movers 10.7 and 11.10 attend to the same name as the name movers and write
\emph{against} it, and the backup name movers become active when the name movers are removed. We saw the trace
of both in \cref{fig:lens,tab:heads}.

\paragraph{How good is the explanation?} The circuit of \citet{wang2023interpretability} has 26 heads in seven
classes (the count in the arXiv version of the paper and in the released code). On their distribution of
prompts the full model has an average logit difference of $3.56$. With every head outside the circuit
mean-ablated, the difference between circuit and full model is $13\%$ of that value, so the circuit retains
$87\%$. This is a strong result, and it also marks the limits of the best-understood circuit in a real language
model. It concerns one narrow task on templated prompts. A part of the behaviour remains unexplained. Several
classes of heads (negative and backup name movers) are described without a clear account of why training
produced them. And the role of the MLPs is left open.

\begin{caution}
Our own figures on 12 prompts reproduce the observational and single-node signatures of this circuit. They are
not a validation of it. We did not run path patching, and we did not test faithfulness, completeness or
minimality.
\end{caution}

\subsection{A network that was understood completely}
\label{sec:modadd}

For one family of small networks the reverse engineering has been carried to the end.
\citet{nanda2023progress} trained a one-layer transformer to add two numbers modulo $P=113$. The input is the
sequence $a$, $b$, \tok{=} with $a,b\in\{0,\dots,P-1\}$, and the correct output is the number
$c\in\{0,\dots,P-1\}$ with $c\equiv a+b \pmod P$. The network was trained on $30\%$ of all pairs. It first
memorises these and much later, quite suddenly, becomes correct on the rest. This delayed generalisation is
called \emph{grokking} \citep{power2022grokking}.

One would expect the network to learn something like a look-up table. It does something else. It uses a small
set of integer frequencies $r$ (five in the main model of the paper) and proceeds as follows.
\begin{enumerate}
\item The embedding maps the input $a$ to the numbers $\cos(2\pi ra/P)$ and $\sin(2\pi ra/P)$ for each of these
frequencies, and likewise for $b$.
\item Attention and MLP multiply these numbers and combine the products, by the addition theorems, into the
cosine and sine of the sum. For the cosine,
\[
\cos\!\Big(\frac{2\pi r(a+b)}{P}\Big)=\cos\!\Big(\frac{2\pi ra}{P}\Big)\cos\!\Big(\frac{2\pi rb}{P}\Big)
-\sin\!\Big(\frac{2\pi ra}{P}\Big)\sin\!\Big(\frac{2\pi rb}{P}\Big).
\]
\item The unembedding multiplies these with $\cos(2\pi rc/P)$ and $\sin(2\pi rc/P)$ for every candidate $c$.
By the addition theorem once more,
\[
\cos\!\Big(\frac{2\pi r(a+b)}{P}\Big)\cos\!\Big(\frac{2\pi rc}{P}\Big)
+\sin\!\Big(\frac{2\pi r(a+b)}{P}\Big)\sin\!\Big(\frac{2\pi rc}{P}\Big)
=\cos\!\Big(\frac{2\pi r(a+b-c)}{P}\Big).
\]
The logit of $c$ is approximately proportional to the sum of these terms over the frequencies,
\begin{equation}\label{eq:modadd}
\sum_{r}\cos\!\Big(\frac{2\pi r\,(a+b-c)}{P}\Big).
\end{equation}
\end{enumerate}

Why does \eqref{eq:modadd} single out the correct answer? If $c\equiv a+b$, every cosine equals $1$. For any
other $c$ the cosines of different frequencies point in different directions and largely cancel. With all
frequencies the cancellation is exact.

\begin{proposition}[Cosines over all frequencies form an indicator]\label{prop:cos}
Let $P\ge2$ be an integer and let $q$ be an integer. Then
\[
\sum_{r=0}^{P-1}\cos\!\Big(\frac{2\pi rq}{P}\Big)=\begin{cases}P&\text{if $q$ is divisible by $P$,}\\ 0&\text{otherwise.}\end{cases}
\]
\end{proposition}
\begin{proof}
If $q$ is divisible by $P$, every angle $2\pi rq/P$ is a multiple of $2\pi$, every term is $1$ and the sum is
$P$. Now let $q$ not be divisible by $P$. Then $\pi q/P$ is not a multiple of $\pi$, so $\sin(\pi q/P)\ne0$.
We multiply the sum by $2\sin(\pi q/P)$ and use the identity $2\sin\phi\cos\psi=\sin(\psi+\phi)-\sin(\psi-\phi)$
with $\phi=\pi q/P$ and $\psi=2\pi rq/P$.
\begin{align*}
2\sin\!\Big(\frac{\pi q}{P}\Big)\sum_{r=0}^{P-1}\cos\!\Big(\frac{2\pi rq}{P}\Big)
&= \sum_{r=0}^{P-1}\Big[\sin\!\Big(\frac{(2r+1)\pi q}{P}\Big)-\sin\!\Big(\frac{(2r-1)\pi q}{P}\Big)\Big] && \by{by the identity}\\
&= \sin\!\Big(\frac{(2P-1)\pi q}{P}\Big)-\sin\!\Big(\frac{-\pi q}{P}\Big) && \by{telescoping sum}\\
&= \sin\!\Big(2\pi q-\frac{\pi q}{P}\Big)+\sin\!\Big(\frac{\pi q}{P}\Big) && \by{since $\sin$ is odd}\\
&= -\sin\!\Big(\frac{\pi q}{P}\Big)+\sin\!\Big(\frac{\pi q}{P}\Big) \;=\; 0. && \by{period $2\pi$, $q$ an integer}
\end{align*}
Dividing by $2\sin(\pi q/P)\ne0$ gives the claim.
\end{proof}

With $q=a+b-c$, the proposition says that the sum \eqref{eq:modadd} over \emph{all} frequencies is $P$ for the
correct $c$ and $0$ for every other $c$. The trained network uses only five frequencies, and this is enough.
Since $P$ is prime, $\cos(2\pi rq/P)=1$ with $0<r<P$ forces $q$ to be divisible by $P$. So the sum of five
cosines equals $5$ at the correct answer and is strictly smaller everywhere else. \Cref{fig:modadd} shows how
large the margin is.

\begin{figure}[t]
\centering
\begin{tikzpicture}
\begin{axis}[width=0.86\textwidth, height=4.6cm, xmin=-58, xmax=58, ymin=-4.3, ymax=5.6, axis lines=left,
  xlabel={$a+b-c$}, ylabel={sum of five cosines}, label style={font=\small}, tick label style={font=\footnotesize},
  xtick={-56,-28,0,28,56}, ytick={-4,-2,0,2,4}, ymajorgrids, grid style={black!10}]
\addplot[ycomb, tutblue, line width=0.7pt, mark=*, mark size=0.9pt, samples at={-56,...,56}]
  {cos(360*14*x/113)+cos(360*35*x/113)+cos(360*41*x/113)+cos(360*42*x/113)+cos(360*52*x/113)};
\end{axis}
\end{tikzpicture}
\caption{The sum \eqref{eq:modadd} for $P=113$ and the five frequencies $r\in\{14,35,41,42,52\}$ reported by
\citet{nanda2023progress}, as a function of $a+b-c$. The value is $5$ at the correct answer $a+b-c=0$ and at
most $2.59$ at every other residue.}
\label{fig:modadd}
\end{figure}

The authors did not guess this algorithm. They found periodic structure in the embedding matrix, traced it
through the network and confirmed each step by interventions. Removing the key frequencies destroys the
performance, and removing everything else improves it. They report that \eqref{eq:modadd} accounts for $95\%$ of
the variance of the logits.

Understanding the mechanism paid off in a second way. Knowing which weights belong to the trigonometric
algorithm, the authors could measure its strength during training. The algorithm does not appear at the moment
when the test loss falls. It is built up gradually while the network still relies on memorisation, and the
sudden improvement happens when weight decay finally removes the memorised solution. The three phases are
called memorisation, circuit formation and cleanup. A curve that looks like a sudden jump from the outside is a
gradual process on the inside.

\subsection{Scaling up: automatic discovery and feature circuits}
\label{sec:featurecircuits}

The IOI circuit took months of manual work on one task in a small model. Two lines of research try to scale
the approach.

\paragraph{Automatic search.} ACDC \citep{conmy2023towards} starts from the full graph and visits its edges
from the output towards the input. An edge is removed if patching it changes the output distribution by less
than a threshold. What remains is the circuit. Edge attribution patching \citep{syed2024attribution} replaces
the exact patching by the linear estimate of \cref{sec:attrpatching}, which makes the search fast. The
numbers that such methods report depend on the kind of ablation and on other details of the procedure
\citep{miller2024transformer}, and should be compared only under identical settings.

\paragraph{Better units.} Heads and neurons are the units that the architecture provides. As
\cref{sec:features} explained, they are mixtures of many features. \emph{Feature circuits} use SAE latents as
nodes \citep{marks2025sparse}. The most developed form is the \emph{attribution graph} of
\citet{ameisen2025circuit}. Its construction has four steps.
\begin{enumerate}
\item All MLPs are replaced by transcoders (\cref{sec:transcoders}), in a variant that lets a latent write to
several later layers. \citet{lindsey2025biology} use 30 million latents for the model Claude 3.5 Haiku.
\item The analysis is done for one prompt. The attention patterns and the scales in $\LN$ are frozen at the
values they take on this prompt.
\item At every layer and position an \emph{error node} is added, equal to the difference between the true MLP
output and the transcoder output. The resulting \emph{local replacement model} then reproduces the activations
of the real model on this prompt exactly.
\item In the local replacement model, every latent influences every later latent through linear maps only. The
direct effect of one latent on another is given by a formula, and the latents and effects that matter for the
output form the attribution graph.
\end{enumerate}
The graph is a hypothesis about the real model. It is tested by interventions in the real model: one suppresses
or activates a latent and checks that the output changes as the graph predicts.

Two findings of \citet{lindsey2025biology} show what such graphs can reveal. For the prompt asking for the
capital of the state that contains Dallas, the graph contains latents for Texas between the input and the answer
\tok{Austin}: the model performs the two steps internally. When writing rhyming verse, the model activates
candidate rhyme words at the end of a line \emph{before} it writes the next line, and then composes the line
so that it ends in one of them. The model plans ahead.

\begin{caution}
An attribution graph describes one prompt. It takes the attention patterns as given, so it does not explain
why the model attends where it does. The error nodes can carry a substantial share of the effect, and that
share is unexplained by construction. And the latents inherit every limitation of \cref{sec:sae}. The honest
summary of a graph is ``this is a mechanism consistent with the model's computation on this prompt, as far as
the replacement captures it''.
\end{caution}

\begin{takeaways}
A circuit is a subgraph of components with a description of what each one computes, and it has to be tested
for faithfulness, completeness and minimality. In small models circuits can be read from the weights. The
induction circuit and the IOI circuit show how heads compose into algorithms, and the modular addition network
shows that a learned algorithm can be understood completely and can differ from anything a programmer would
write. Attribution graphs extend the approach to large models, one prompt at a time and with an unexplained
remainder.
\end{takeaways}

\begin{exercise}\label{ex:skiptrigram}
In the model \eqref{eq:onelayer} with a single head, suppose the head makes ``\tok{keep} $\dots$ \tok{in}
$\to$ \tok{mind}'' and ``\tok{keep} $\dots$ \tok{at} $\to$ \tok{bay}'' likely. Using the three tables of
\cref{sec:pathexp}, argue that it will also raise the probability of ``\tok{keep} $\dots$ \tok{in} $\to$
\tok{bay}'' to some extent. Which table would have to depend on the destination token to avoid this?
\end{exercise}
\begin{exercise}\label{ex:induction}
Explain why the induction circuit of \cref{sec:induction} cannot be implemented by a transformer with one
layer. Use the fact that in \eqref{eq:path} the attention score depends only on the tokens $t_i$ and $t_j$ and
their positions.
\end{exercise}
\begin{exercise}\label{ex:cosprime}
Let $P=113$ and $r=14$. For which integers $q$ is $\cos(2\pi rq/P)=1$? Why does the argument fail for $P=112$?
\end{exercise}
